\begin{abstract}
Current Vision-Language-Action (VLA) models are often constrained by a rigid, static interaction paradigm, which lacks the ability to see, hear, speak, and act concurrently as well as handle real-time user interruptions dynamically. This hinders seamless human-robot collaboration, resulting in an inflexible and unresponsive user experience. To address these limitations, we introduce VITA-E, a novel human-robot interaction framework designed for both behavioral concurrency and nearly real-time interruption. The core of our approach is a dual-model architecture where two parallel VLA instances operate as an ``Active Model'' and a ``Standby Model'', allowing the robot to observe its environment, listen to user speech, provide verbal responses, and execute actions, all concurrently and interruptibly, mimicking human-like multitasking capabilities. We further propose a ``model-as-controller'' paradigm, where we fine-tune the VLM to generate special tokens that serve as direct system-level commands, coupling the model's reasoning with the system's behavior. Experiments conducted on a physical humanoid robot demonstrate that VITA-E can reliably handle complex interactive scenarios. Our framework is compatible with various dual-system VLA models, achieving an extremely high success rate on emergency stops and speech interruptions while also successfully performing concurrent speech and action. This represents a significant step towards building more natural and capable robotic assistants. 
% Our homepage is \url{https://lxysl.github.io/VITA-E/}.


\end{abstract}